% Lines 3261-3812 from original attention_transformers_notes_ghost.tex
% Chapter 12: Advanced Training Techniques

\section{Advanced Training Techniques}
%==============================================================================

\subsection{Optimization Strategies for Transformers}

\subsubsection{Learning Rate Scheduling}

The original transformer paper introduced a custom learning rate schedule:

\begin{equation}
\text{lr} = d_{\text{model}}^{-0.5} \cdot \min(\text{step}^{-0.5}, \text{step} \cdot \text{warmup}^{-1.5})
\end{equation}

\begin{notationbox}
\textbf{Notation:}
\begin{itemize}
    \item $d_{\text{model}}$: Model dimension
    \item $\text{step}$: Current training step
    \item $\text{warmup}$: Number of warmup steps
\end{itemize}
\end{notationbox}

Let's implement various learning rate schedules:

\begin{lstlisting}[language=Python, caption=Advanced Learning Rate Schedules]
from __future__ import annotations
import math
from abc import ABC, abstractmethod
from dataclasses import dataclass
import torch.optim as optim

@dataclass(frozen=True, slots=True)
class ScheduleConfig:
    optimizer: optim.Optimizer
    d_model: int
    warmup_steps: int = 4000

class LRScheduler(ABC):
    """Abstract base class for learning rate schedulers"""

    def __init__(self, config: ScheduleConfig) -> None:
        self.config = config
        self.step_num = 0

    @abstractmethod
    def get_lr(self) -> float:
        """Get current learning rate"""
        pass

    def step(self) -> None:
        """Update learning rate"""
        self.step_num += 1
        lr = self.get_lr()
        for param_group in self.config.optimizer.param_groups:
            param_group['lr'] = lr

class TransformerLRScheduler(LRScheduler):
    """Original transformer learning rate schedule"""

    def get_lr(self) -> float:
        step = max(1, self.step_num)
        scale = self.config.d_model ** -0.5

        return scale * min(
            step ** -0.5,
            step * self.config.warmup_steps ** -1.5
        )

class CosineAnnealingWarmup(LRScheduler):
    """Cosine annealing with linear warmup"""

    def __init__(
        self,
        config: ScheduleConfig,
        max_lr: float,
        min_lr: float,
        total_steps: int
    ) -> None:
        super().__init__(config)
        self.max_lr = max_lr
        self.min_lr = min_lr
        self.total_steps = total_steps

    def get_lr(self) -> float:
        if self.step_num < self.config.warmup_steps:
            # Linear warmup
            return self.max_lr * self.step_num / self.config.warmup_steps
        else:
            # Cosine annealing
            progress = (self.step_num - self.config.warmup_steps) / \
                      (self.total_steps - self.config.warmup_steps)
            return self.min_lr + (self.max_lr - self.min_lr) * \
                   0.5 * (1 + math.cos(math.pi * progress))

class PolynomialDecayWarmup(LRScheduler):
    """Polynomial decay with warmup"""

    def __init__(
        self,
        config: ScheduleConfig,
        max_lr: float,
        end_lr: float,
        power: float,
        total_steps: int
    ) -> None:
        super().__init__(config)
        self.max_lr = max_lr
        self.end_lr = end_lr
        self.power = power
        self.total_steps = total_steps

    def get_lr(self) -> float:
        if self.step_num < self.config.warmup_steps:
            # Linear warmup
            return self.max_lr * self.step_num / self.config.warmup_steps
        elif self.step_num >= self.total_steps:
            return self.end_lr
        else:
            # Polynomial decay
            decay_steps = self.total_steps - self.config.warmup_steps
            decay_progress = (self.step_num - self.config.warmup_steps) / decay_steps
            return self.end_lr + (self.max_lr - self.end_lr) * \
                   (1 - decay_progress) ** self.power
\end{lstlisting}

\subsubsection{Gradient Accumulation and Mixed Precision}

Handle large models with limited memory:

\begin{lstlisting}[language=Python, caption=Gradient Accumulation with Mixed Precision]
from torch.cuda.amp import autocast, GradScaler

class TransformerTrainer:
    """Advanced trainer for transformer models"""

    def __init__(
        self,
        model: nn.Module,
        optimizer: optim.Optimizer,
        accumulation_steps: int = 4,
        use_amp: bool = True,
        clip_grad_norm: float = 1.0
    ) -> None:
        self.model = model
        self.optimizer = optimizer
        self.accumulation_steps = accumulation_steps
        self.use_amp = use_amp
        self.clip_grad_norm = clip_grad_norm

        # Mixed precision training
        self.scaler = GradScaler() if use_amp else None

        # Training statistics
        self.step = 0
        self.accumulated_loss = 0.0

    def train_step(
        self,
        batch: dict[str, Tensor],
        criterion: nn.Module
    ) -> dict[str, float]:
        """Single training step with gradient accumulation"""

        # Adjust loss for accumulation
        loss_scale = 1.0 / self.accumulation_steps

        # Forward pass with mixed precision
        if self.use_amp:
            with autocast():
                outputs = self.model(
                    batch['input_ids'],
                    attention_mask=batch.get('attention_mask')
                )
                loss = criterion(outputs, batch['labels'])
                loss = loss * loss_scale
        else:
            outputs = self.model(
                batch['input_ids'],
                attention_mask=batch.get('attention_mask')
            )
            loss = criterion(outputs, batch['labels'])
            loss = loss * loss_scale

        # Backward pass
        if self.use_amp:
            self.scaler.scale(loss).backward()
        else:
            loss.backward()

        # Accumulate loss for logging
        self.accumulated_loss += loss.item()

        # Update weights every accumulation_steps
        if (self.step + 1) % self.accumulation_steps == 0:
            # Gradient clipping
            if self.use_amp:
                self.scaler.unscale_(self.optimizer)

            torch.nn.utils.clip_grad_norm_(
                self.model.parameters(),
                self.clip_grad_norm
            )

            # Optimizer step
            if self.use_amp:
                self.scaler.step(self.optimizer)
                self.scaler.update()
            else:
                self.optimizer.step()

            self.optimizer.zero_grad()

            # Log accumulated loss
            avg_loss = self.accumulated_loss
            self.accumulated_loss = 0.0

            self.step += 1

            return {'loss': avg_loss, 'step': self.step}

        self.step += 1
        return {}
\end{lstlisting}

\subsubsection{Dynamic Batching and Sequence Packing}

Optimize training efficiency with variable-length sequences:

\begin{lstlisting}[language=Python, caption=Dynamic Batching Implementation]
from typing import Sequence
import numpy as np

@dataclass(frozen=True)
class DynamicBatch:
    """Batch with sequences of varying lengths"""
    input_ids: Tensor
    attention_mask: Tensor
    labels: Tensor | None
    seq_lens: list[int]

class DynamicBatcher:
    """Create efficient batches from variable-length sequences"""

    def __init__(
        self,
        max_tokens: int,
        max_sequences: int,
        pad_token_id: int
    ) -> None:
        self.max_tokens = max_tokens
        self.max_sequences = max_sequences
        self.pad_token_id = pad_token_id

    def create_batches(
        self,
        sequences: list[list[int]],
        labels: list[list[int]] | None = None
    ) -> list[DynamicBatch]:
        """Create batches optimizing for GPU utilization"""

        # Sort sequences by length for better packing
        sorted_indices = np.argsort([len(seq) for seq in sequences])[::-1]
        sorted_sequences = [sequences[i] for i in sorted_indices]
        sorted_labels = [labels[i] for i in sorted_indices] if labels else None

        batches = []
        current_batch = []
        current_labels = [] if labels else None
        current_tokens = 0

        for i, seq in enumerate(sorted_sequences):
            seq_len = len(seq)

            # Check if adding this sequence exceeds limits
            if current_batch:
                max_len_in_batch = max(len(s) for s in current_batch)
                new_max_len = max(max_len_in_batch, seq_len)
                new_total_tokens = new_max_len * (len(current_batch) + 1)

                if (new_total_tokens > self.max_tokens or
                    len(current_batch) >= self.max_sequences):
                    # Finalize current batch
                    batch = self._create_batch(
                        current_batch,
                        current_labels
                    )
                    batches.append(batch)

                    # Start new batch
                    current_batch = [seq]
                    current_labels = [sorted_labels[i]] if labels else None
                else:
                    current_batch.append(seq)
                    if labels:
                        current_labels.append(sorted_labels[i])
            else:
                current_batch = [seq]
                current_labels = [sorted_labels[i]] if labels else None

        # Don't forget last batch
        if current_batch:
            batch = self._create_batch(current_batch, current_labels)
            batches.append(batch)

        return batches

    def _create_batch(
        self,
        sequences: list[list[int]],
        labels: list[list[int]] | None
    ) -> DynamicBatch:
        """Create padded batch from sequences"""

        seq_lens = [len(seq) for seq in sequences]
        max_len = max(seq_lens)
        batch_size = len(sequences)

        # Pad sequences
        padded_input = torch.full(
            (batch_size, max_len),
            self.pad_token_id,
            dtype=torch.long
        )

        attention_mask = torch.zeros(
            (batch_size, max_len),
            dtype=torch.float
        )

        for i, seq in enumerate(sequences):
            padded_input[i, :len(seq)] = torch.tensor(seq)
            attention_mask[i, :len(seq)] = 1.0

        # Handle labels if provided
        padded_labels = None
        if labels:
            padded_labels = torch.full(
                (batch_size, max_len),
                -100,  # Ignore index for loss
                dtype=torch.long
            )
            for i, label_seq in enumerate(labels):
                padded_labels[i, :len(label_seq)] = torch.tensor(label_seq)

        return DynamicBatch(
            input_ids=padded_input,
            attention_mask=attention_mask,
            labels=padded_labels,
            seq_lens=seq_lens
        )
\end{lstlisting}

\subsection{Regularization Techniques}

\subsubsection{Attention Dropout Patterns}

Different dropout strategies for attention:

\begin{lstlisting}[language=Python, caption=Advanced Attention Dropout]
class StructuredAttentionDropout(nn.Module):
    """Structured dropout for attention matrices"""

    def __init__(
        self,
        drop_prob: float,
        drop_mode: str = "whole_rows"  # "whole_rows", "whole_cols", "blocks"
    ) -> None:
        super().__init__()
        self.drop_prob = drop_prob
        self.drop_mode = drop_mode

    def forward(
        self,
        attention_weights: Tensor,
        training: bool = True
    ) -> Tensor:
        if not training or self.drop_prob == 0:
            return attention_weights

        batch_size, n_heads, seq_len, seq_len = attention_weights.shape

        if self.drop_mode == "whole_rows":
            # Drop entire rows (queries)
            mask = torch.rand(batch_size, n_heads, seq_len, 1) > self.drop_prob
            mask = mask.float().to(attention_weights.device)
            return attention_weights * mask / (1 - self.drop_prob)

        elif self.drop_mode == "whole_cols":
            # Drop entire columns (keys)
            mask = torch.rand(batch_size, n_heads, 1, seq_len) > self.drop_prob
            mask = mask.float().to(attention_weights.device)
            return attention_weights * mask / (1 - self.drop_prob)

        elif self.drop_mode == "blocks":
            # Drop contiguous blocks
            block_size = 4
            n_blocks = seq_len // block_size

            # Create block mask
            block_mask = torch.rand(
                batch_size, n_heads, n_blocks, n_blocks
            ) > self.drop_prob

            # Expand to full size
            mask = block_mask.repeat_interleave(
                block_size, dim=2
            ).repeat_interleave(
                block_size, dim=3
            )

            # Handle remainder
            if seq_len % block_size != 0:
                full_mask = torch.ones(
                    batch_size, n_heads, seq_len, seq_len
                )
                full_mask[:, :, :mask.shape[2], :mask.shape[3]] = mask
                mask = full_mask

            mask = mask.float().to(attention_weights.device)
            return attention_weights * mask / (1 - self.drop_prob)

class AdaptiveDropout(nn.Module):
    """Dropout rate that adapts during training"""

    def __init__(
        self,
        initial_drop: float,
        final_drop: float,
        warmup_steps: int,
        decay_steps: int
    ) -> None:
        super().__init__()
        self.initial_drop = initial_drop
        self.final_drop = final_drop
        self.warmup_steps = warmup_steps
        self.decay_steps = decay_steps
        self.current_step = 0

    def get_drop_prob(self) -> float:
        """Calculate current dropout probability"""
        if self.current_step < self.warmup_steps:
            # Linear warmup
            return self.initial_drop * self.current_step / self.warmup_steps
        elif self.current_step < self.warmup_steps + self.decay_steps:
            # Linear decay
            progress = (self.current_step - self.warmup_steps) / self.decay_steps
            return self.initial_drop + (self.final_drop - self.initial_drop) * progress
        else:
            return self.final_drop

    def forward(self, x: Tensor) -> Tensor:
        drop_prob = self.get_drop_prob()
        return F.dropout(x, p=drop_prob, training=self.training)

    def step(self) -> None:
        """Update internal step counter"""
        self.current_step += 1
\end{lstlisting}

\subsubsection{Stochastic Depth for Transformers}

Apply stochastic depth (layer dropout) to transformer blocks:

\begin{lstlisting}[language=Python, caption=Stochastic Depth Implementation]
class StochasticDepthTransformerLayer(nn.Module):
    """Transformer layer with stochastic depth"""

    def __init__(
        self,
        d_model: int,
        n_heads: int,
        d_ff: int,
        drop_path_rate: float = 0.1
    ) -> None:
        super().__init__()

        self.attention = nn.MultiheadAttention(
            d_model, n_heads, batch_first=True
        )
        self.ff = nn.Sequential(
            nn.Linear(d_model, d_ff),
            nn.GELU(),
            nn.Linear(d_ff, d_model)
        )

        self.norm1 = nn.LayerNorm(d_model)
        self.norm2 = nn.LayerNorm(d_model)

        self.drop_path_rate = drop_path_rate

    def forward(
        self,
        x: Tensor,
        training: bool = True
    ) -> Tensor:
        # Stochastic depth - skip entire layer with probability
        if training and torch.rand(1).item() < self.drop_path_rate:
            return x

        # Standard transformer layer
        attn_out, _ = self.attention(x, x, x)
        x = x + self._drop_path(attn_out, training)
        x = self.norm1(x)

        ff_out = self.ff(x)
        x = x + self._drop_path(ff_out, training)
        x = self.norm2(x)

        return x

    def _drop_path(
        self,
        x: Tensor,
        training: bool
    ) -> Tensor:
        """Apply drop path if not already skipped"""
        if not training:
            return x

        # Scale by survival probability
        return x / (1 - self.drop_path_rate)

class ProgressiveStochasticDepth(nn.Module):
    """Transformer with progressive stochastic depth"""

    def __init__(
        self,
        n_layers: int,
        d_model: int,
        n_heads: int,
        d_ff: int,
        max_drop_rate: float = 0.3
    ) -> None:
        super().__init__()

        # Linearly increase drop rate
        drop_rates = torch.linspace(0, max_drop_rate, n_layers)

        self.layers = nn.ModuleList([
            StochasticDepthTransformerLayer(
                d_model, n_heads, d_ff, drop_rate.item()
            )
            for drop_rate in drop_rates
        ])

    def forward(self, x: Tensor) -> Tensor:
        for layer in self.layers:
            x = layer(x, training=self.training)
        return x
\end{lstlisting}

%==============================================================================
